\section{Exact eight-arm construction predicates}
\label{eight:specification}
This section specifies the predicates $\mathsf A_8^\epsilon$, $\mathsf O_8^\epsilon$ and $\mathsf H_8^\epsilon$ used in the article, with $\epsilon=o$ or $e$. Within each of the two domains below we suppress the subscript and superscript: $\A$, $\mathsf O$ and $\Hh$ always mean that domain's eight-arm predicates. They do not mean the ten-family predicate $\mathsf A_7$. Set $K(I)=\sum_{i\in I}\floor{L_i/2}$; Outer-route complements are taken in the eight-element index set; within a recursive rule, a complementary pair or triple is taken within the selected subset. All offsets are nonnegative integers.

\begin{lemma}\label{eight:single-equal-offsets}
An arm of length $L$ admits every offset $0\le k<L$.  More generally,
for $r,L\ge1$, the equal-arm spider $\Sp{L^r}$ admits every such offset.
\end{lemma}
\begin{proof}
For the single arm, take an ordinary graceful permutation of
$1,\ldots,L$ beginning at $L-k$, and add $k$ to each entry.  The root
edge has color $L$, and the internal colors are $1,\ldots,L-1$.
The prescribed-first-label path theorem is
Lemma~\ref{common:prescribed-first-path}; it includes the ordinary
non-alpha permutations needed at the exceptional odd-order requests.
For the equal-arm assertion, subdivide the root-zero $r$-leaf unit star
by factor $L$.  Theorem~\ref{common:offset-subdivision}, applied at source
offset zero, gives every residue $0,\ldots,L-1$.
\end{proof}

\subsection{The twelve odd-pair clauses}\label{eight:odd-pairs}

For $a=2t+1\le b=2u+1$, where $t,u\ge1$, define $P_o(a,b;k)$ to mean
$k\ge0$ and at least one of the following twelve conditions.  Parameters
introduced in one condition are local to that condition.

\begin{enumerate}
\item[(P1)] $u\le k\le a-1$, and either
  $4k\ne4a-b-3$ or $3\mid a$ and $3\mid b$.
\item[(P2)] $k\le t$, and either $4k\ne a-1$ or $2b\le3a-1$.
\item[(P3)] $k\le u$, $k+u\le a-1$, and $4k\ne b-1$.
\item[(P4)] $a=b$ and $k<a$.
\item[(P5)] $a=b\ge5$ and $2k\le3a-5$.
\item[(P6)] $a\ge7$, $b=a+2$, and $k=2a$.
\item[(P7)] $a\ge7$, $b=a+2$, $a\le k$, and $2k\le3a-1$.
\item[(P8)] With $S=k-2u-2$, $\delta=u-t$, and $V=k-t-u+1$,
  \[
   S\ge2,\quad t\ge2S+2,\quad 0\le\delta\le S-1,
   \quad\delta\ne S-3,\quad1\le V\le t.
  \]
\item[(P9)] With $S=k-2t$ and $V=k-t-u+1$,
  \[
   1\le S\le t\le u\le2S-2,\quad u>S,
   \quad u\ne2t-S,\quad1\le V\le u.
  \]
\item[(P10)] With $S=k-2t$, $g=u-2S$, $h=k-t-u$, and $V=k-t-u+1$,
  \[
   g\ge1,\quad h\ge2,\quad S\ge g+h+1,\quad1\le V\le u.
  \]
\item[(P11)] With $S=u-t+1$ and $g=2t-u-2$,
  \[
   k=t+u+2,\quad g\ge1,\quad S\ge g+3,\quad u\ge3.
  \]
\item[(P12)] With $S=3u-k+3$, $h=t+2u-k+1$,
  $g=2k-4u-t-4$, and $V=k-t-u+1$,
  \[
   g\ge1,\quad h\ge2,\quad S\ge g+h+1,\quad1\le V\le t.
  \]
\end{enumerate}

\begin{lemma}\label{eight:odd-pair-sufficiency}
If $P_o(a,b;k)$ holds, then $k\in\Om(\Sp{a,b})$.
\end{lemma}
\begin{proof}
Clauses (P1)--(P3) are the HIGH, LOW, and reversed-LOW constructions,
with their exact exceptions and repairs, proved in
Section~\ref{app:pair-offsets}.  Clause (P4) is
Lemma~\ref{eight:single-equal-offsets}.  The extended equal-odd interval
(P5), the consecutive interval (P7), the two general kernel bands
(P8)--(P9), and the three asymmetric bands (P10)--(P12) are constructed
in Section~\ref{construction:odd-pair-kernels}.  The consecutive maximum
(P6) also follows from the complete maximal-pair theorem in
Section~\ref{app:pair-offsets}.

In (P12) the kernel construction interchanges the order of the two
arms.  This does not change the rooted spider, but the restriction
$V\le t$ is part of its even-path premise and must be retained.
The consecutive construction includes its fifteen literal small kernels;
the kernel-band proofs include all source endpoint and hole conditions.

The divisibility repair in (P1) is used only within its source scope.
Indeed, if $a=3$, HIGH gives $u\le k\le2$.  If also $3\mid b$ and
$b\ge a$, this forces $b=3$.  Its exceptional equality would be
$4k=6$, impossible for integer $k$.  Every active repaired exceptional
request therefore has $a,b\ge9$, so its quotient lengths are at least
three.  No small forbidden source request is used.
\end{proof}

\subsection{Finite recursive construction rules}\label{eight:recursions}

This section first treats odd lengths $L_i\ge3$.  Every index subset
inherits the increasing order of the full list.  The same recursions
will later be used for one even minimum, with the pair predicate and
the insertion capacities changed as explicitly stated there.

\begin{definition}\label{eight:alpha-rule}
The predicate $\A(I)$ is defined by the following finite recursion.
\begin{enumerate}
\item For $|I|=1$ or $2$, it is true.
\item For $|I|=3$, it means $\max_{i\in I}L_i\ge5$.
\item For $|I|\ge4$, it holds if either $|I|=4$ and all four lengths
  are equal odd integers at least seven, or some $j\in I$ satisfies
  \[
   \A(I\setminus\{j\}),\qquad
   L_j>2K(I\setminus\{j\}),\qquad
   L_j\ne4K(I\setminus\{j\})+1.
  \]
\end{enumerate}
\end{definition}

These are root-zero alpha constructions of index $K(I)$.  The triple
rule uses the complete odd triple theorem: among odd lengths at least
three, the only omitted profile is $(3,3,3)$
(Section~\ref{construction:all-odd-alpha3}).  The equal-four rule uses
Section~\ref{construction:equal-odd-alpha4}; its lower bound seven is
retained.  The last rule is Lemma~\ref{common:alpha-extension}, with
both its strict threshold and its exceptional length intact.

\begin{definition}\label{eight:offset-rule}
Require $k\ge0$.  For one arm, $\mathsf O(\{i\},k)$ means $k<L_i$;
for two arms, it means $P_o(L_i,L_j;k)$.  For three ordered lengths
$x\le y\le z$, let $T(x,y,z;k)$ be the disjunction of:
\begin{enumerate}
\item $k=0$;
\item $y=z$, $\floor{y/2}\le k\le x-1$, and either $x=y$ or both
  \[
   4k\ne3(y-1),\qquad 4k\ne4x-y-3;
  \]
\item $x=z$ and $k<x$;
\item some arm $L_j$ satisfies
  \[
   k\le\floor{L_j/2},\qquad4k\ne L_j-1,
  \]
  and its complementary pair has predicate $P_o$ at offset
  $k+\floor{L_j/2}$.
\end{enumerate}
This is $\mathsf O(I,k)$ when $|I|=3$.  For $|I|=4$, let
$\mathsf O(I,k)$ hold if all four lengths agree and $k$ is below that
common length, or if some $j\in I$ satisfies the same two displayed
conditions and its complementary triple satisfies $T$ at offset
$k+\floor{L_j/2}$.  Put $\mathsf O(I,k)=\mathrm{false}$ for $|I|\ge5$.
\end{definition}

The triple rule at $k=0$ uses ordinary hub-zero gracefulness, not an
alpha assertion about $(3,3,3)$.  The repeated-long rule and the one-leg
insertion rule, including their excluded equalities, are proved in
Section~\ref{app:triple-offsets}.  The equal-arm rules follow from
Lemma~\ref{eight:single-equal-offsets}.  Thus each true offset predicate
has the stated offset realization.

\begin{definition}\label{eight:hub-rule}
Put $\Hh(I)=\mathrm{true}$ for $|I|\le3$.  For four ordered lengths
$a\le b\le c\le d$, let $\Hh(I)$ be the disjunction of:
\begin{enumerate}
\item $a\le5$;
\item $a=b$ and $c=d$;
\item some triple $J\subset I$ has $\A(J)$, and its complementary
  arm $j$ satisfies $L_j>K(J)$;
\item some arm $j$ has
  $\mathsf O(I\setminus\{j\},\floor{L_j/2})$;
\item some pair $J\subset I$ has $\mathsf O(I\setminus J,K(J))$.
\end{enumerate}
For $|I|\ge5$, put
\[
 \Hh(I)=\A(I)\ \lor\!
 \bigvee_{\varnothing\ne J\subsetneq I}
 \bigl(\A(J)\land\mathsf O(I\setminus J,K(J))\bigr).
\]
\end{definition}

\begin{lemma}\label{eight:predicate-soundness}
Every true $\A(I)$ supplies a root-zero alpha labeling of index $K(I)$;
every true $\mathsf O(I,k)$ supplies offset $k$; and every true $\Hh(I)$
supplies an ordinary hub-zero graceful labeling.
\end{lemma}
\begin{proof}
The alpha and offset assertions follow by induction on arity from the
construction lemmas just specified.  The first two four-arm hub rules
follow from Lemma~\ref{seven:four-short} and
Section~\ref{po:paired-four}, respectively. The other three are
alpha-offset amalgamation, using respectively an offset single arm,
an alpha single arm, or an alpha pair.  The pair-root exception $(2,2)$
cannot occur on odd lengths.  The higher-arity hub rule is exactly
alpha-offset amalgamation, or the alpha labeling itself.  Its
singleton-offset term is the ordinary arm-attachment condition $L>K$;
it must not be confused with the alpha-preserving condition
$L>2K$, $L\ne4K+1$.
\end{proof}

\subsection{One even minimum}\label{eight:even-minimum}

For this section put
\[
 L_0=2h_0,\qquad L_i=2h_i+1\quad(1\le i\le7),
 \qquad2\le L_0<L_1\le\cdots\le L_7,
 \qquad L_1\ge3.
\]
The even length is the unique even arm and the shortest arm of this
eight-arm core.  The equality $L_0=L_1$ is impossible by parity.

Every one- or two-arm selection is root-zero alpha, since $(2,2)$
cannot occur.  For three arms the exact sufficient predicate remains
\begin{equation}\label{eight:small-triple-guard}
 \A(I)\quad\Longleftrightarrow\quad\max_{i\in I}L_i\ge5.
\end{equation}
Here this is a statement about the specified sufficient predicate.
A triple avoiding index zero is odd and its only omitted profile is
$(3,3,3)$.  A triple containing zero has EOO parity with an even
minimum, and its only omitted profile is $(2,3,3)$.
The complete positive constructions, including minima two and four,
are in Sections~\ref{construction:all-odd-alpha3},
\ref{construction:all-eoo-alpha3}, and~\ref{app:small-alpha}.
Thus no unconditional all-triples alpha shortcut is used.
For larger alpha pieces keep Definition~\ref{eight:alpha-rule}; its
equal-four clause is available only on four odd arms of common length
at least seven.

\subsubsection{The mixed-pair predicate}\label{eight:mixed-pairs}

For any selected ordered pair $a\le b$, put
$t=\floor{a/2}$, $u=\floor{b/2}$ and
$c(a)=\ceil{a/2}-1$, $c(b)=\ceil{b/2}-1$.  Define $P_e(a,b;k)$ to
require $k\ge0$ and one of the following six base clauses:
\begin{enumerate}
\item[(M1)] $u\le k\le a-1$, and either $4k\ne4a-b-3$ or
  $a\ge6$, $3\mid a$, and $3\mid b$.
\item[(M2)] $k\le c(a)$, and either $4k\ne a-1$ or $2b\le3a-1$.
\item[(M3)] $k\le c(b)$, $k+u\le a-1$, and $4k\ne b-1$.
\item[(M4)] $a=b$ and $k<a$.
\item[(M5)] $k=a+b-2$, with $(a,b)\notin\{(2,2),(3,3)\}$.
\item[(M6)] $3\mid a$, $3\mid b$, $(a,b)\notin\{(6,6),(9,9)\}$,
  and $a+b-6\le k\le a+b-4$.
\end{enumerate}
If both selected arms are odd, also allow exactly clauses
(P5)--(P12), with their displayed guards.  None of those eight clauses
is applied to a mixed pair.

\begin{lemma}\label{eight:mixed-pair-sufficiency}
Each clause defining $P_e(a,b;k)$ supplies an offset-$k$ labeling.
\end{lemma}
\begin{proof}
Clauses (M1)--(M3) are the mixed HIGH, LOW and reversed-LOW
constructions in Section~\ref{app:pair-offsets}.  The extra guard
$a\ge6$ in the divisible-by-three repair gives shortest quotient at
least two; it is not removed at the present minimum.  Clause (M4)
is Lemma~\ref{eight:single-equal-offsets}.  Clause (M5) is the complete
maximal-pair theorem, including both small exceptions, in
Section~\ref{app:pair-offsets}.  For (M6), apply that theorem to
$(a/3,b/3)$ and then apply factor-three offset subdivision.  The
three residues give the interval
\[
 3\bigl((a+b)/3-2\bigr)+\{0,1,2\}
   =[a+b-6,a+b-4].
\]
The two excluded target pairs are exactly the forbidden quotient pairs
$(2,2)$ and $(3,3)$.  Quotient $(1,1)$ is allowed; thus this clause
also validly supplies offsets $0,1,2$ on $(3,3)$.  It does not supply
the forbidden maximal offset four there.  The remaining odd-only
clauses follow from Lemma~\ref{eight:odd-pair-sufficiency}.
\end{proof}

Use the same offset and hub recursions as before, replacing $P_o$ by
$P_e$ and replacing each one-leg insertion capacity
$\floor{L_j/2}$ by
\[
 c(L_j)=\ceil{L_j/2}-1.
\]
The shifted source offset is still $k+\floor{L_j/2}$, and the guard
$4k\ne L_j-1$ is retained.  Equal-arm and repeated-long direct clauses
are unchanged.  The ordinary $k=0$ triple rule refers to hub-zero
gracefulness, so it makes no alpha claim about either exceptional
triple.  In the four-arm hub rule the alpha-triple test
\eqref{eight:small-triple-guard} remains explicit.

